% !TEX program = xelatex
% ============================================================
% 符号清单（工作记录文件，不进论文正文）
% 用途：随正文撰写逐节登记已出现的字母记号及含义，统一口径。
% 论文正文中不设符号表（含义在首次出现处行内解释）。
% ============================================================
\documentclass[11pt]{ctexart}
\usepackage[a4paper,margin=2.4cm]{geometry}
\usepackage{amsmath,amssymb,booktabs}
% 记号宏（与 main.tex 导言区一致）
\newcommand{\E}{\mathbb{E}}
\newcommand{\Var}{\operatorname{Var}}

\begin{document}

\begin{center}
  {\LARGE\bfseries 符号清单（工作记录）}
\end{center}
\vspace{0.5em}

本文正文不设符号表，各记号含义在首次出现处以行内文字解释。本文件随撰写进度逐节登记，供后续各节写作时统一口径、避免记号冲突。

\section*{II.A 含背景光的直方图观测模型（\texttt{2\_method/2\_1.tex}）}

\begin{center}
\begin{tabular}{@{}llp{0.58\textwidth}@{}}
\toprule
符号 & 含义 & 备注 \\
\midrule
$p$ & 目标像素（像元）索引 & \\
$b$ & 直方图的离散 bin 坐标 & 每 bin 宽 $66.9\,\mathrm{ps}$，对应 $1\,\mathrm{cm}$ 深度 \\
$H_p(b)$ & 像素 $p$ 的直方图在 bin $b$ 的光子计数 & 随机变量 \\
$\lambda_p(b;\mu_p)$ & $H_p(b)$ 的期望（泊松参数） & $N_p h(b-\mu_p)+\beta_p$ \\
$\mu_p$ & 像素 $p$ 的真实回波中心坐标 & 待估参数 \\
$N_p$ & 像素 $p$ 的期望信号光子数 & \\
$h(\cdot)$ & 归一化系统脉冲响应（IRF） & $\sum_b h(b-\mu_p)=1$ \\
$\beta_p$ & 每 bin 平均背景光子与暗计数 & 平坦背景假设下为常数 \\
$\beta_p(b)$ & 随 bin 变化的背景计数 & 背景非平坦时的推广形式 \\
$\E[\cdot]$ & 期望算子 & \\
$\Var[\cdot]$ & 方差算子 & \\
$\mathcal I_p(\mu_p)$ & 像素 $p$ 含背景泊松模型下 $\mu_p$ 的 Fisher 信息 & \\
$h'(\cdot)$ & IRF 的导数 & \\
$\hat\mu_p$ & 由直方图构造的回波中心估计量 & Cram\'er--Rao 下界针对其方差 \\
$v_p$ & 回波中心估计方差 & $v_p\approx\mathcal I_p^{-1}$；高斯无背景时 $\sigma_h^2/N_p$ \\
$\sigma_h^2$ & 高斯 IRF 的方差 & \\
\bottomrule
\end{tabular}
\end{center}

\section*{II.B 死时间畸变与 Coates 校正（\texttt{2\_method/2\_2.tex}）}

\begin{center}
\begin{tabular}{@{}llp{0.58\textwidth}@{}}
\toprule
符号 & 含义 & 备注 \\
\midrule
$r_p(b)$ & 单个测距周期内 bin $b$ 的期望光子数（到达率） & $\lambda_p(b;\mu_p)=M\,r_p(b)$ \\
$R_p(b)$ & 累积率 & $R_p(b)=\sum_{b'\le b}r_p(b')$ \\
$M$ & 每帧有效测距周期数 & 实验中 $M=1125$ \\
$H_p^{\mathrm{raw}}(b)$ & 原始（含死时间畸变）直方图计数 & \\
$\mathrm{cum}(b)$ & 原始直方图累积计数 & $\mathrm{cum}(b)=\sum_{b'\le b}H_p^{\mathrm{raw}}(b')$ \\
$\hat r_p(b)$ & Coates 校正后的到达率估计 & 式：$\ln\frac{M-\mathrm{cum}(b-1)}{M-\mathrm{cum}(b)}$ \\
$F_p(b)$ & 校正后的率域直方图 & $F_p(b)\triangleq\hat r_p(b)$；后续对齐、匹配与融合的输入 \\
$\mathcal I_p^{\mathrm{eff}}$ & 校正后直方图的有效 Fisher 信息 & $v_p^{\mathrm{eff}}=1/\mathcal I_p^{\mathrm{eff}}$；稀疏区 $=\mathcal I_p$ \\
$A_p(b)$ & 逐 bin 放大（惩罚）因子 & $e^{R_p(b-1)}(e^{r_p(b)}-1)/r_p(b)\ge1$；稀疏区 $\approx1$ \\
\bottomrule
\end{tabular}
\end{center}

\section*{II.C 单目相对深度引导的局部对齐（\texttt{2\_method/2\_3.tex}）}

\begin{center}
\begin{tabular}{@{}llp{0.58\textwidth}@{}}
\toprule
符号 & 含义 & 备注 \\
\midrule
$\mathcal{G}_p$ & 像素 $p$ 的邻域（邻居集合） & 原记 $\mathcal{N}_p$，为避免与 $N_p$ 及高斯分布 $\mathcal{N}$ 混淆改名 \\
$b_p^M$ & 像素 $p$ 的单目深度坐标（bin 域） & 单目相对深度经逐场景最小二乘标定映射到 bin 域 \\
$e_p$ & 单目预测误差 & $b_p^M=\mu_p+e_p$ \\
$\Delta_{pq}$ & 单目相对位移 & $\Delta_{pq}=b_p^M-b_q^M=(\mu_p-\mu_q)+(e_p-e_q)$ \\
$k_{pq}^{M}$ & 粗对齐量（搜索中心） & $k_{pq}^{M}=\operatorname{round}(\Delta_{pq})$ \\
$\rho_{pq}(\delta)$ & 归一化（循环）互相关匹配得分 & $\in[0,1]$；在 $F_p$ 与平移后的 $F_q$ 之间计算 \\
$\widetilde F_{pq}(b)$ & 对齐到目标坐标系的邻居直方图 & $\widetilde F_{pq}(b)=F_q(b-\hat k_{pq})$；$\hat k_{pq}=k_{pq}^{M}+\hat\delta_{pq}$ \\
\bottomrule
\end{tabular}
\end{center}

\section*{II.D 可融合性概率模型（\texttt{2\_method/2\_4.tex}）}

\begin{center}
\begin{tabular}{@{}llp{0.58\textwidth}@{}}
\toprule
符号 & 含义 & 备注 \\
\midrule
$z_{pq}$ & 二值潜变量：$q$ 与 $p$ 深度相近、可融合 & $z_{pq}=1$ 可融合，$z_{pq}=0$ 不可融合 \\
$\sigma_s$ & 可融合状态下 $\Delta_{pq}$ 的散布 & 并入表面内真实深度起伏与单目误差之差 \\
$\sigma_d$ & 不可融合状态下 $\Delta_{pq}$ 的散布 & 含跨表面深度跳变；$\sigma_d>\sigma_s$ \\
$\Lambda_{\Delta}$ & 几何似然比 & 随 $|\Delta_{pq}|$ 单调递减的软门控 \\
$\Lambda_{\rho}$ & 匹配证据似然比 & 幂式随 $\rho$ 单调递增（$a_s>a_d$、$b_s<b_d$ 时） \\
$a_s,b_s,a_d,b_d$ & 条件 Beta 分布参数 & 可依赖 $\Delta_{pq}$；由标注数据拟合 \\
$\pi_{pq}$ & 先验可融合概率 & $\pi_{pq}=\mathbb{P}(z_{pq}=1)$ \\
$\gamma_{pq}$ & 后验可融合概率 & $\gamma_{pq}=\mathbb{P}(z_{pq}=1\mid\Delta_{pq},\rho_{pq})$ \\
\bottomrule
\end{tabular}
\end{center}

\end{document}
